% ============================================================================
%  CAPÍTULO I: INTRODUCCIÓN
% ============================================================================

\section{Introducción}

En toda organización, la infraestructura de red constituye el pilar sobre el
cual se sustentan las operaciones diarias y la comunicación interna. Dentro de
esta infraestructura, los servidores de correo electrónico y los servidores de
archivos desempeñan un rol fundamental: el primero permite el intercambio formal
de información entre los miembros de la organización, mientras que el segundo
centraliza el almacenamiento y la distribución de documentos\footnote{Stewart
y Chapple, 2017.}

El protocolo \textit{Simple Mail Transfer Protocol} (SMTP) es el estándar
utilizado globalmente para la transferencia de correo electrónico entre servidores
y para el envío de mensajes desde clientes hacia el servidor\footnote{Postel,
1982.} En combinación con protocolos de acceso como \textit{Internet Message
Access Protocol} (IMAP) y \textit{Post Office Protocol} (POP3), se completa el
ciclo de vida de un correo: desde su recepción en el servidor hasta su lectura
por parte del usuario final.

Por otro lado, el protocolo \textit{Server Message Block} (SMB) —implementado
en Linux mediante Samba— permite compartir archivos y carpetas entre máquinas
de distintos sistemas operativos en una red local, ofreciendo permisos de
acceso a nivel de recurso compartido y a nivel NTFS
(Linux equivalent)\footnote{Tanenbaum y Wetherall, 2017.}

El presente informe documenta la configuración de un servidor de correo
electrónico y la planificación de un servidor de archivos dentro del proyecto
integrador del curso de Taller de Sistemas Operativos II. Para el correo se
implementó un entorno basado en Postfix (MTA) y Dovecot (servidor IMAP/POP3)
sobre Debian, con resolución de nombres vía BIND9 y autenticación de usuarios
a través del subsistema SASL integrado en Dovecot. El servidor de archivos
utiliza Samba (SMB) con carpetas compartidas y permisos de acceso por usuario,
siendo desarrollado por otro integrante del equipo.

Este capítulo introduce el contexto general del taller y presenta la importancia
de estos servicios de red en la administración de sistemas informáticos.
